\chapter{Comment s'entraîner}\label{ch:iv1:how-to-train}

Ce chapitre indique ce que testent les entretiens quantitatifs, dans quel ordre ils le testent, et comment utiliser ces trois livres pour se préparer. On le lit une fois avant de commencer, et on y revient quand on planifie une semaine. Il ne contient aucune question.

Ce qui suit sur les sociétés provient de deux types de sources~: ce que les sociétés publient sur leur propre processus, qui relève surtout de la logistique, et ce que les candidats rapportent publiquement, qui est détaillé mais non audité. Il faut donc considérer chaque chiffre ici comme une indication plutôt que comme une spécification. Aucune société n'est nommée, car un processus varie selon le bureau, le poste et l'année, et un nom suggérerait une précision que personne ne possède.

\section{Les étapes}

La plupart des processus suivent le même schéma, avec des étapes allant de la moins coûteuse et automatique à la plus coûteuse et personnelle.

\textbf{Candidature.} Un CV, parfois accompagné d'un relevé de notes ou d'une courte lettre de motivation. Certaines sociétés invitent à un test en ligne chaque candidat dont le CV passe un filtre rapide. Un grand teneur de marché indique que son processus commence par une candidature et des tests en ligne, suivis d'un entretien comportemental avec un recruteur.

\textbf{Test en ligne.} Une ou plusieurs épreuves chronométrées, passées seul, avant qu'un être humain ne s'intéresse de près au candidat. C'est à ce stade, selon les témoignages, que se terminent la plupart des candidatures, bien qu'aucune société ne publie cette proportion. La \Cref{sec:iv1:how-to-train:tests} décrit leur contenu.

\textbf{Présélection enregistrée ou avec un recruteur.} Un bref appel, ou une vidéo à laquelle on répond face à un écran, sans personne de l'autre côté. La page destinée aux étudiants d'une grande banque décrit un entretien vidéo d'environ 30 minutes. Les candidats rapportent un mélange de questions de motivation et de calcul mental léger, qui sert en partie à vérifier la manière dont le candidat réagit.

\textbf{Tours techniques.} Un à trois appels de 30 à 60 minutes, par téléphone ou en visioconférence, avec une personne qui occupe le poste. C'est le cœur de ce livre et des deux suivants.

\textbf{Travail à domicile.} Ce n'est pas universel, mais il est signalé surtout dans les fonds dont les chercheurs travaillent sur des données~: un exercice sur données d'une durée allant de quelques heures à deux jours, suivi d'un appel destiné à vérifier que l'on comprend son propre travail. Pour les analystes de recherche en investissement fondamental, l'équivalent signalé est un stock pitch et une étude de cas écrite avec un modèle financier, avec environ une semaine pour la réaliser.

\textbf{Dernier tour.} Plusieurs entretiens successifs dans une même journée, souvent appelés superday. La page destinée aux étudiants d'une grande banque en indique deux à cinq pour les recrutements sur le campus, selon la division. Les témoignages concernant les ingénieurs le situent plutôt autour de quatre à six. Il faut s'attendre à des versions plus approfondies des tours techniques, à quelques questions d'adéquation au profil et, dans de nombreuses sociétés de trading, à un jeu de market making.

\textbf{Décision.} Selon les témoignages, la décision tombe d'un jour ou deux à une semaine après le dernier tour. Un processus pour un stagiaire ou un jeune diplômé dure, selon les témoignages, de trois à huit semaines au total.

\begin{dated}{2026-11}{Les tentatives sont limitées}\label{dat:iv1:how-to-train:retakes}
Un grand teneur de marché indique sur sa page de recrutement publiée que ses tests en ligne ne peuvent être passés qu'une fois tous les huit mois, qu'aucune réinitialisation n'est possible une fois une tentative commencée, et que les résultats récents sont valables pour un poste similaire. Vérifier cela pour chaque société. Ne pas passer un test réel pour s'entraîner, et ne pas en commencer un avec une connexion médiocre. Faire ses exercices sur les feuilles de ce livre.
\end{dated}

\section{Tests en ligne}\label{sec:iv1:how-to-train:tests}

Les candidats signalent quatre ingrédients, dans des proportions variables~: le calcul mental, les probabilités et l'espérance, les régularités et la logique, et la programmation. Voici les formats qu'ils décrivent le plus souvent. Tous proviennent de témoignages de candidats recueillis par des sites de préparation, et aucun n'est publié par une société.

\begin{dated}{2026-11}{Formats signalés par les candidats}\label{dat:iv1:how-to-train:formats}
\textbf{Calcul contre la montre.} Le format le plus connu consiste en 80 questions à traiter en 8 minutes, sans calculatrice. Des variantes de 60 questions en 6 minutes et de 50 en 5 minutes sont aussi signalées, certaines avec quatre options par question et d'autres avec des réponses saisies. Les témoignages divergent sur la pénalité~: la plupart indiquent qu'une mauvaise réponse fait perdre un point, et au moins un parle de deux. Les candidats signalent aussi une note de passage, toujours non officielle, d'environ 55 sur la version à 80 questions.

\textbf{Suites.} Trouver le terme suivant, avec cinq options. Dans un cas signalé, 26 questions en 25 minutes~; dans un autre, 30 questions en environ 6 minutes.

\textbf{Probabilités et résolution de problèmes.} Les exemples signalés sont 15 à 17 questions en une heure, environ 6 en 45 minutes, et 18 en 30 minutes. Les questions prennent d'autant plus de temps que leur nombre diminue.

\textbf{Jeux cognitifs et questionnaires.} De courts jeux qui mesurent l'attention et la mémoire de travail, parfois accompagnés d'un questionnaire de personnalité.

\textbf{Programmation.} Pour les développeurs et de nombreux chercheurs, en général deux ou trois problèmes en 60 à 120 minutes sur une plateforme hébergée, notés automatiquement à l'aide de cas de test cachés.
\end{dated}

Les feuilles d'exercices de vitesse du \cref{ch:iv1:speed-drills} varient en longueur, en durée et en barème, afin de ne pas s'habituer à un seul format. Avant de passer un test réel, il faut se renseigner sur son barème et établir la règle à appliquer pour deviner. Une réponse devinée, qui vaut $+a$ si elle est juste et $-b$ si elle est fausse, n'est avantageuse que si la probabilité d'avoir juste dépasse $b/(a + b)$. Cette seule ligne détermine le nombre de questions à laisser sans réponse.

\section{Ce qu'on demande à chaque rôle}

Les étiquettes de chaque question indiquent à quel profil elle correspond. Voici la tendance générale qui les sous-tend. Les témoignages varient, et les intitulés varient encore davantage~: il faut donc lire la description du poste, et non son intitulé.

\textbf{Trader et trader quantitatif.} Calcul mental, probabilités sans calculatrice, jeux de paris et de market making. Dans les sociétés qui forment leurs propres traders, on attend des candidats qu'ils réfléchissent vite et à voix haute, et un grand teneur de marché indique clairement qu'il n'évalue pas les candidats au trading sur la finance ni sur l'économie. Les desks des banques et des fonds attendent une meilleure connaissance des marchés. Commencer par les parties de ce livre sur la vitesse et les probabilités, puis passer aux jeux de trading du livre~19.

\textbf{Chercheur quantitatif.} Probabilités et statistiques en profondeur, régression et apprentissage automatique, un peu de programmation, et discussion de ses propres projets. Les candidats en doctorat sont souvent invités à expliquer leurs recherches. Les parties de ce livre consacrées aux probabilités et à la statistique viennent en premier, puis les études de cas de recherche du livre~19 et le traitement des données du livre~20.

\textbf{Analyste.} Ici, le livre désigne à la fois les analystes orientés données et les analystes fondamentaux. Selon les témoignages, un analyste fondamental dans un fonds doit présenter un stock pitch, en position longue et courte, dans presque tous les entretiens, ainsi qu'une valorisation et un modèle. Le calcul et les estimations de ce livre aident, mais l'essentiel du travail se trouve dans les chapitres sur la valorisation et sur le pitch du livre~19.

\textbf{Data scientist.} Statistiques, régression, apprentissage automatique, raisonnement expérimental et traitement des données, avec de la programmation en Python et en SQL. Il faut s'appuyer sur la partie statistique de ce livre et sur les chapitres consacrés aux données du livre~20.

\textbf{Développeur quantitatif et développeur logiciel.} Algorithmes et structures de données, maîtrise approfondie d'un langage, un peu de probabilités et un peu de conception de systèmes. Le livre~20 est l'ouvrage principal, et les échauffements en probabilités de ce livre suffisent pour l'essentiel des probabilités rencontrées.

\textbf{Ingénieur basse latence.} Les témoignages décrivent un ou deux tests de programmation en direct d'environ 60 minutes, suivis d'une journée de quatre à six tours portant sur le fonctionnement interne du C++, la concurrence, le comportement du cache et la conception de systèmes, et parfois un bref tour quantitatif. Le livre~20 est l'ouvrage principal.

\textbf{Quant de risque et quant de banque.} Les témoignages décrivent des probabilités et de la statistique, suivies du pricing, du calcul stochastique, de C++ ou de Python et, pour les postes de validation, de la manière de revoir le modèle d'autrui. Il faut s'appuyer sur les chapitres de ce livre consacrés à la statistique et au calcul stochastique, puis sur les chapitres sur le pricing et les risques du livre~19.

Le type de société modifie aussi cette répartition. Les teneurs de marché penchent vers la vitesse et les jeux, les fonds vers la profondeur de la recherche et les travaux propres du candidat, les banques vers la connaissance des produits, le calcul stochastique et l'adéquation au profil, et les sociétés crypto vers les mêmes mathématiques, auxquelles s'ajoute la structure de marché propre à ce secteur. Ce ne sont que des tendances, et l'on rencontre des exceptions à chaque tour.

\section{Comment s'entraîner}

Il faut utiliser les trois livres de façon active. La suite de cette section expose la méthode du livre, et trois études la justifient.

\textbf{Répondre avant de lire.} S'engager sur une réponse, sur papier ou à voix haute, avant de regarder les réponses en fin d'ouvrage. Dans une expérience bien connue, les étudiants qui s'étaient testés sur un passage en ont retenu davantage une semaine plus tard que ceux qui l'avaient relu, bien que ces derniers se sentent plus sûrs d'eux. Lire une réponse sans avoir cherché procure cette même fausse assurance.

\textbf{Espacer son entraînement.} Une synthèse des travaux sur l'espacement montre que répartir l'entraînement dans le temps est plus efficace que le bachotage, et que le meilleur intervalle croît avec la durée pendant laquelle on doit retenir la matière. Il faut donc consacrer vingt minutes par jour pendant six semaines plutôt qu'un long week-end. Revenir sur une question manquée après quelques jours, puis après environ deux semaines.

\textbf{Parler.} La plupart des tours se déroulent à l'oral. Il faut dire ce que l'on fait, énoncer ses hypothèses, donner tôt une réponse approximative et l'affiner. Le recruteur ne peut aider que s'il entend où en est le candidat, et il note la façon de raisonner autant que le point d'arrivée. Si l'on s'entraîne seul, il faut s'enregistrer en train de répondre à voix haute à une question à une étoile, puis l'écouter une fois.

\textbf{Vérifier chaque réponse avant de l'énoncer.} Tester un petit cas, contrôler les ordres de grandeur et regarder les unités, et utiliser un cas limite. Des contrôles rapides pour le calcul sont donnés au \cref{ch:iv1:mental-technique}. Une erreur que l'on repère soi-même joue en sa faveur, et celle que l'on annonce avec assurance ne joue pas.

\textbf{S'entraîner dans les conditions de l'épreuve.} Une méta-analyse sur la répétition des tests montre que les scores aux tests cognitifs augmentent d'environ un quart d'écart-type lors d'une seconde passation, grâce à la pratique et au coaching. L'entraînement aide donc sur le format, et il vaut mieux avoir déjà tiré profit de cet effet avant la tentative réelle. Il faut utiliser les feuilles d'exercices et les épreuves chronométrées avec un minuteur, d'une seule traite, et se corriger honnêtement.

\section{Un plan}

Les étoiles donnent un temps cible (voir la préface), et les tableaux de la préface comptent les questions par chapitre pour chaque rôle. Un ordre raisonnable pour la plupart des candidats~:

\begin{enumerate}
\item \textbf{Mesurer.} Passer une feuille d'exercices et la première épreuve chronométrée sans préparation. Cela indique sur quoi travailler, et évite de passer des semaines sur ce que l'on maîtrise déjà.
\item \textbf{La vitesse, un peu chaque jour.} Dix minutes par jour sur une feuille d'exercices ou une famille d'exercices, pendant toute la préparation. La vitesse demande une pratique régulière, et un peu chaque jour ne coûte presque rien.
\item \textbf{Notions essentielles.} Travailler les chapitres qu'exige le poste, dans l'ordre à l'intérieur de chaque chapitre. Arrêter un chapitre quand les questions à trois étoiles se passent bien, et sauter ce qu'excluent les étiquettes. Une question résolue dans le temps indiqué par les étoiles est considérée comme acquise. Une question manquée est inscrite sur une liste à reprendre.
\item \textbf{Épreuves mélangées.} À partir du milieu de la préparation, passer une épreuve chronométrée par semaine. Des questions classées par thème dévoilent la méthode. Les épreuves chronométrées ne le font pas~: elles sont la seule mesure honnête pour savoir si l'on possède la méthode ou si l'on se contente de la reconnaître.
\item \textbf{La dernière ligne droite.} Pendant les deux dernières semaines, renoncer aux nouveaux contenus. Reprendre la liste, passer les épreuves chronométrées restantes et s'exercer à parler en déroulant un entretien blanc du livre~19.
\end{enumerate}

Si l'on n'a qu'une semaine, il faut faire une feuille d'exercices par jour, les dix premières questions de chaque chapitre de probabilités et une épreuve chronométrée, et laisser le reste de côté. Un livre de cette taille est conçu pour une préparation complète, et il n'est pas nécessaire de le terminer. Il vaut mieux en faire la moitié avec soin que tout en gardant les réponses ouvertes.
